\scanpage{323}
\mspage{239}
% archive sheet number: 322

l'image de $x_0$ dans $I$. Quant à $\mathrm{PF}_{\unsure{\pi|\kappa}}2$, si on a des \diagram{\begin{tikzcd}[ampersand replacement=\&,cramped,column sep=small] x_0 \arrow[r,shift left=1.2ex,"u_\alpha"] \arrow[r] \arrow[r,shift right=1.2ex] \& x_1 \end{tikzcd}} dans $J$, d'où des \deleted{\unsure{$g(x_1)$|$g(x_0)$}} \diagram{\begin{tikzcd}[ampersand replacement=\&,cramped,column sep=small] g(x_0) \arrow[r,shift left=1.2ex,"g(u_{\unsure{1|\alpha|2}})"] \arrow[r,no head] \arrow[r,shift right=1.2ex] \& g(x_1) \end{tikzcd}}, on égalise \inserted{dans $I$} par $g(x_1)\xrightarrow{w}y$, mais comme $g$ est coétal, $\exists$ $x_1\xrightarrow{v}x_2$ dans $J$ telle que $g(v)=w$, donc on a $g(vu_\alpha)=g(\unsure{w|v}u_\beta)$ pour $\alpha,\beta\in A$, et comme $g$ est fidèle, on en conclut $vu_\alpha=vu_\beta$, \deleted{qed.} donc $v$ égalise $(u_\alpha)_{\alpha\in A}$, qed.

\deleted{\uline{Corollaire} 4.9.6.9.3. Soit $I$ une catégorie, $\pi$ un cardinal infini. Pour que $I$ soit \unsure{grde|gde} devant $\pi$, il f. et s. que \deleted{\unsure{$I$ soit $0$-connexe|$I$ soit connexe}} pour tout $i\in\mathop{\mathrm{Ob}}\nolimits I$, $ {}_i\backslash I$ le soit, \uline{et} que $I$ soit $0$-connexe.}

% slanted
\margin{Faux déjà dans le cas filtrant sans plus, cf \unsure{ex.|Ex.} page 236}

\deleted{C'est nécessaire par 4.9.6.9.1, puisque les $ {}_i\backslash I$ sont $0$-connexes et coétales sur $I$. C'est suffisant, car \deleted{\unsure{les|le cas}} cette condition implique que les $ {}_i\backslash I$ satisfont $\mathrm{PF}_\pi 1$, et par \deleted{\unsure{la|le}} 4.9.6.9.2 il \deleted{suffit} reste à vérifier \uline{$\mathrm{PF}_\pi 2$ pour $I$}.}

Prouvons $b_\pi)\longrightarrow a_\pi)$ : Si $I'$ \deleted{\unsure{est|est grde}} et les $ {}_{i'}\backslash I$ sont \unsure{grdes|gdes} devant $\pi$, alors $I$ l'est aussi (et $f$ est cofinal, les $ {}_{i'}\backslash I$ étant $0$-connexes, car filtrantes). Comme on sait déjà que $I$ est $0$-connexe ($I'$ l'étant et $f$ cofinal), il reste, par 4.9.6.9.2, à vérifier $\mathrm{PF}_\pi 1$ et $\mathrm{PF}_\pi 2$.

